\documentclass[tikz,border=8pt]{standalone}
\usepackage{tikz}
\usepackage{amsmath}
\usepackage{amssymb}
\usepackage{pgfplots}
\pgfplotsset{compat=1.18}
\usepackage{ctex}
\usetikzlibrary{calc,positioning,arrows.meta,trees}

% LC 102 层序遍历 BFS
% 树：[3,9,20,null,null,15,7]
% BFS 队列演化：[3] -> [9,20] -> [15,7] -> []

\begin{document}
\begin{tikzpicture}[
  every node/.style={},
  level distance=1.4cm,
  level 1/.style={sibling distance=3.6cm},
  level 2/.style={sibling distance=2.0cm},
  edge from parent/.style={draw,->,>=Stealth,thick}
]

%% ── 标题 ──────────────────────────────────────────────────
\node[font=\large\bfseries] at (2.5, 1.6) {LC 102 二叉树层序遍历（BFS）};
\node[font=\small,gray] at (2.5, 1.1) {树 \texttt{[3,9,20,null,null,15,7]}，按层收集节点};

%% ── 主树 ────────────────────────────────────────────────
\node[draw,circle,minimum size=0.78cm,fill=blue!10,font=\bfseries] (root) at (2.5,0) {3}
  child {
    node[draw,circle,minimum size=0.78cm,fill=blue!10,font=\bfseries] (n9) {9}
  }
  child {
    node[draw,circle,minimum size=0.78cm,fill=blue!10,font=\bfseries] (n20) {20}
    child {
      node[draw,circle,minimum size=0.78cm,fill=blue!10,font=\bfseries] (n15) {15}
    }
    child {
      node[draw,circle,minimum size=0.78cm,fill=blue!10,font=\bfseries] (n7) {7}
    }
  };

%% ── 层标注（右侧括号） ────────────────────────────────────
\node[font=\small,gray] at (5.6,  0.0) {$\leftarrow$ Level 0};
\node[font=\small,gray] at (5.6, -1.4) {$\leftarrow$ Level 1};
\node[font=\small,gray] at (5.6, -2.8) {$\leftarrow$ Level 2};

%% ── BFS 队列快照（下方 4 个） ────────────────────────────
\node[font=\small\bfseries] at (2.5, -4.2) {BFS 队列演化快照};

% 快照框样式：draw=gray!60, fill=white, rounded corners
% 快照 1
\node[draw=gray!60,fill=white,rounded corners=3pt,inner sep=6pt,
      font=\small,align=center,minimum width=2.8cm]
  at (-1.0, -5.0)
  {\textbf{步骤 1}\\队列：\texttt{[3]}\\出队：3\\入队：9, 20};

% 快照 2
\node[draw=gray!60,fill=white,rounded corners=3pt,inner sep=6pt,
      font=\small,align=center,minimum width=2.8cm]
  at (2.0, -5.0)
  {\textbf{步骤 2}\\队列：\texttt{[9,20]}\\出队：9, 20\\入队：15, 7};

% 快照 3
\node[draw=gray!60,fill=white,rounded corners=3pt,inner sep=6pt,
      font=\small,align=center,minimum width=2.8cm]
  at (5.0, -5.0)
  {\textbf{步骤 3}\\队列：\texttt{[15,7]}\\出队：15, 7\\入队：（无）};

% 快照 4
\node[draw=gray!60,fill=white,rounded corners=3pt,inner sep=6pt,
      font=\small,align=center,minimum width=2.8cm]
  at (8.0, -5.0)
  {\textbf{步骤 4}\\队列：\texttt{[]}\\\textit{遍历完成}};

%% ── 箭头连接快照 ─────────────────────────────────────────
\draw[->,>=Stealth,gray!70,thick]  (0.5,-5.0) -- (0.9,-5.0);
\draw[->,>=Stealth,gray!70,thick]  (3.5,-5.0) -- (3.9,-5.0);
\draw[->,>=Stealth,gray!70,thick]  (6.5,-5.0) -- (6.9,-5.0);

%% ── 最终结果 ─────────────────────────────────────────────
\node[draw=gray!50,fill=green!5,thin,rounded corners=3pt,
      font=\small,inner sep=6pt,align=center]
  at (3.5, -6.8)
  {\textbf{输出（levels）}：$[[3],\ [9,20],\ [15,7]]$};

\end{tikzpicture}
\end{document}
